\documentclass[a4paper,12pt]{article}   % 文档类：A4纸张，正文字体大小为12pt

\usepackage[UTF8]{ctex}                 % 支持中文排版（UTF-8编码）
\usepackage{xcolor}                     % 颜色支持（用于定义和调用颜色）
\usepackage{minted}                     % 代码高亮宏包（依赖 Pygments）
\usepackage{fontspec}                   % 字体选择宏包（用于设置系统字体，如等宽字体）

% 字体
\setmainfont{TeX Gyre Termes}
\setmonofont{Cascadia Code}
\setCJKmainfont{Noto Serif CJK SC}

% 代码风格
\usemintedstyle{gruvbox-light}      % 主题
\definecolor{codebg}{HTML}{F3F3F3}  % 背景色

% 全局设置
\setminted{
    % linenos,                % 显示代码行号
    breaklines,             % 自动换行（超长行折行）
    frame=none,             % 不绘制代码框边框
    framesep=2mm,           % 代码内容与边框（或页面边缘）的间距
    fontsize=\small,        % 代码字体大小（small 约为 9pt）
    bgcolor=codebg,         % 代码块背景色（使用自定义颜色 codebg）
    tabsize=4,              % 制表符占用的空格数
    autogobble,             % 自动去除代码左侧多余的公共空白（缩进）
    xleftmargin=0.5em       % 代码块左侧额外缩进距离
}

\title{minted 代码高亮示例}
\author{}
\date{}

\begin{document}

\maketitle

\section{Python}

\begin{minted}{python}
def fibonacci(n):
    if n <= 1:
        return n
    return fibonacci(n - 1) + fibonacci(n - 2)

print([fibonacci(i) for i in range(10)])
\end{minted}

\section{C++}

\begin{minted}{cpp}
#include <iostream>

int main() {
    std::cout << "Hello C++" << std::endl;
    return 0;
}
\end{minted}

\section{Rust}

\begin{minted}{rust}
fn main() {
    println!("Hello Rust!");
}
\end{minted}

\section{Java}

\begin{minted}{java}
public class Main {
    public static void main(String[] args) {
        System.out.println("Hello Java");
    }
}
\end{minted}

\section{Shell Script}

\begin{minted}{bash}
#!/bin/bash

echo "Hello Shell"

for f in *.txt; do
    echo "$f"
done
\end{minted}

\section{HTML}

\begin{minted}{html}
<!DOCTYPE html>
<html>
<head>
    <title>Hello</title>
</head>
<body>
    <h1>Hello HTML</h1>
</body>
</html>
\end{minted}

\section{CSS}

\begin{minted}{css}
body {
    font-family: Arial;
    background: #f5f5f5;
}

h1 {
    color: steelblue;
}
\end{minted}

\section{JavaScript}

\begin{minted}{javascript}
function hello(name) {
    console.log(`Hello ${name}`);
}

hello("World");
\end{minted}

\section{JSON}

\begin{minted}{json}
{
    "name": "Alice",
    "age": 20,
    "languages": [
        "Python",
        "Rust",
        "C++"
    ]
}
\end{minted}

\section{TOML}

\begin{minted}{toml}
name = "Alice"
age = 20
languages = ["Python", "Rust", "C++"]
\end{minted}

\section{YAML}

\begin{minted}{yaml}
name: Alice
age: 20
languages:
  - Python
  - Rust
  - C++
\end{minted}

\end{document}